% Generated draft: requires mathematical and source review.
\documentclass[11pt]{article}
\usepackage[margin=25mm]{geometry}
\usepackage{fontspec}
\setmainfont{Latin Modern Roman}
\usepackage{amsmath,amssymb,mathtools,unicode-math}
\setmathfont{Latin Modern Math}
\usepackage{enumitem,needspace,xurl,hyperref}
\setlength{\parindent}{0pt}
\setlength{\parskip}{5pt}
\setlength{\emergencystretch}{3em}
\newcommand{\N}{\mathbb{N}}
\newcommand{\Z}{\mathbb{Z}}
\newcommand{\Q}{\mathbb{Q}}
\newcommand{\R}{\mathbb{R}}
\newcommand{\C}{\mathbb{C}}
\newcommand{\F}{\mathbb{F}}
\newcommand{\A}{\mathbb{A}}
\newcommand{\PP}{\mathbb{P}}
\newcommand{\E}{\mathbb{E}}
\newcommand{\eps}{\varepsilon}
\newcommand{\dd}{\,\mathrm{d}}
\begin{document}
\section*{UnsolvedMath ID: 30006225}
\textbf{Source problem code:} catalogue entry
\subsection{Definitions and mathematical statement}
Let \(S^n=\{\xi\in\R^{n+1}:|\xi|=1\}\) carry its unit round metric \(g_{S^n}\), and write \(\dd\sigma\) for its volume measure. Interpret \(\omega_n=\int_{S^n}\dd\sigma\); this interpretation requires confirmation because the supplied record does not define \(\omega_n\). Write \(L_{2r}\) for the operators appearing in the source and \(Q^{(r)}_g\) for its \(r\)-th \(Q\)-curvature. Their definitions and normalizations are not supplied, and the record does not specify whether the orders \(r\) must be integers.

For the source's parameter range \(0<\ell<k<n/2\), let \(u\in C^\infty(S^n)\) be strictly positive and set
\[
g=u^{4/(n-2k)}g_{S^n},\qquad v=u^{(n-2\ell)/(n-2k)},\qquad \theta=\frac{n-2k}{n-2\ell},\qquad A_r=\frac{\Gamma(n/2+r)}{\Gamma(n/2-r)},
\]
where \(\Gamma\) is the gamma function. The proposed hypothesis is pointwise positivity \(Q^{(k)}_g>0\).

The supplied transcription of inequality (10) has lost its mathematical layout. A reconstruction consistent with its displayed factors and exponents is the conjectural comparison
\[
\int_{S^n}uL_{2k}u\,\dd\sigma
\;\geq\;
A_k A_\ell^{-\theta}\,\omega_n^{\frac{2(k-\ell)}{n-2\ell}}
\left(\int_{S^n}vL_{2\ell}v\,\dd\sigma\right)^\theta.
\]
This formula is a reconstruction requiring comparison with the cited report, not an unqualified recovery of its typesetting. In particular, the missing operator and curvature conventions must be checked before identifying these quadratic expressions with normalized total \(Q\)-curvatures. The real-power expression also requires an appropriate sign for the integral on the right; the supplied record does not explain how this follows from its conventions and hypotheses.

The proposed equality characterization is
\[
u(\xi)=a(1+b\cdot\xi)^{-(n-2k)/2},\qquad \xi\in S^n,\qquad b\in B^{n+1}:=\{b\in\R^{n+1}:|b|<1\}.
\]
The source writes \(a\in\R\). Since \(1+b\cdot\xi>0\) and the admissible \(u\) is strictly positive, the admissible members of this family necessarily have \(a>0\).

The opening source question allows \(0<\ell<k\leq n/2\), whereas its explicit conjecture assumes \(k<n/2\). Only the latter range is formulated here: the displayed conformal factor is undefined at \(k=n/2\), and no critical-order replacement is supplied.

As a supplied partial-result claim, the record reports that for \(n\geq5\), \(k=2\), and \(\ell=1\), the comparison holds under the stated weaker hypothesis that \(g=u^{4/(n-4)}g_{S^n}\) has positive scalar curvature, expressed there as
\[
L_2\!\left(u^{(n-2)/(n-4)}\right)>0.
\]
It reports the same equality family, with exponent \(-(n-4)/2\). This is an attributed result in the supplied record, not an independently verified theorem assertion.
\subsection{Short English statement}
The conjecture proposes a sharp comparison between total \(k\)-th and \(\ell\)-th \(Q\)-curvatures on a conformally round sphere, assuming positive \(k\)-th \(Q\)-curvature, with equality precisely for the displayed conformal factors. The supplied record reports a partial result for \(k=2\), \(\ell=1\), under positive scalar curvature. The general formula, operator conventions, and admissible orders require source verification.
\subsection{Sources}
\begin{itemize}
\item Local UnsolvedMath registry: \texttt{data/unsolvedmath/problems.json}, permanent ID \(30006225\), source problem code \texttt{catalogue entry}, title \emph{Isoperimetric Inequalities Between Total Q-Curvatures}. No individual proposer is named in the supplied record.
\item Supplied report citation: \emph{MATRIX-MFO Tandem Workshop: Nonlinear Geometric Diffusion Equations} (2025), Oberwolfach Reports, DOI \texttt{10.4171/owr/2025/11}, pp.\ 485--530. Supplied source and literature link, labeled \emph{Oberwolfach Report 11/2025}: \url{https://doi.org/10.4171/owr/2025/11}.
\item Literature assessment dated exactly \texttt{2026-08-22}: the supplied curation labels the problem partially solved, reporting the \(k=2\), \(\ell=1\) case under a weaker scalar-curvature hypothesis and describing the general range \(0<\ell<k<n/2\) as conjectural. This is supplied dated curation, not independently verified; its own qualification requests independent verification of time-sensitive status and cited claims.
\item The supplied registry labels its statement \texttt{corrected\_verified} and says that the official report was checked and its notation repaired. These are supplied provenance claims; no independent source visit or mathematical verification is asserted here.
\end{itemize}
\end{document}
